% Bilan des forces sur une personne dans un ascenseur qui démarre vers le haut (L'essentiel, chapitre 14)
\begin{tikzpicture}[x=1cm,y=1cm]
\draw[otInk!70,line width=0.8pt] (-1.4,0) rectangle (1.4,3.6);
\fill[otRule!60] (-1.4,0) rectangle (1.4,0.12);
\filldraw[pcAppareil] (-0.45,0.12) rectangle (0.45,0.3);
% silhouette simplifiée
\fill[otNavy!25] (-0.22,0.3) rectangle (0.22,2.2);
\fill[otNavy!25] (0,2.45) circle (0.22);
\fill[otRed!80!black] (0,1.35) circle (0.05);
\node[pcLegende,anchor=west] at (0.08,1.25) {G};
\draw[pcForce] (0,1.35) -- (0,3.05) node[pcLegende,anchor=west,text=otGreenDark] {$\vec{N}$ : \qty{792}{N}};
\draw[pcForce] (0,1.35) -- (0,-0.15) node[pcLegende,anchor=west,text=otGreenDark] {$\vec{P}$ : \qty{687}{N}};
\draw[pcAcc] (-0.85,1.4) -- (-0.85,2.25) node[pcLegende,anchor=south,text=otRed!80!black] {$\vec{a}$};
\draw[otMuted,line width=0.5pt,-{Stealth[length=1.8mm]}] (2.0,0.5) -- (2.0,1.6) node[pcLegende,anchor=south] {$z$};
\end{tikzpicture}
\par\smallskip\legende{Au démarrage vers le haut, la force exercée par le pèse-personne dépasse le poids : leur somme est dirigée comme $\vec{a}$.}
